% ===== PRÉAMBULE CANONIQUE — à coller VERBATIM en tête de CHAQUE .tex =====
\documentclass[11pt,a4paper]{article}
\usepackage[utf8]{inputenc}\usepackage[T1]{fontenc}\usepackage{lmodern}
\usepackage[french]{babel}
\usepackage{amsmath,amssymb,mathtools}\usepackage{icomma}
\usepackage{siunitx}\sisetup{locale=FR}  % \qty{}{}, \si{}, \ang{}
\usepackage{xcolor}\usepackage{colortbl}
\usepackage{tikz}\usepackage{pgfplots}\pgfplotsset{compat=1.18}
\usepackage[siunitx,europeanresistors]{circuitikz} % circuits (élec.)
\usepackage[most]{tcolorbox}\usepackage{enumitem}\usepackage{array,booktabs}
\usepackage[a4paper,margin=2.2cm]{geometry}
\usetikzlibrary{arrows.meta,calc,angles,quotes,positioning,patterns,%
  backgrounds,shapes.geometric,decorations.pathmorphing,decorations.markings,intersections}
\definecolor{primary}{RGB}{222,49,99}\definecolor{primarydark}{RGB}{150,22,55}
\definecolor{defcol}{RGB}{0,114,178}\definecolor{methcol}{RGB}{52,92,148}
\definecolor{excol}{RGB}{0,135,90}\definecolor{attncol}{RGB}{200,40,55}
\tcbset{boxbase/.style={enhanced,breakable,boxrule=0.9pt,arc=2.5pt,
  left=3mm,right=3mm,top=2.3mm,bottom=2.3mm,fonttitle=\bfseries,coltitle=white}}
% --- boîtes TUTORIEL (noms EXACTS, ne pas renommer) ---
\newtcolorbox{cle}[1][]{boxbase,colback=defcol!6,colframe=defcol,
  title={\textsc{À retenir}\ifx\relax#1\relax\relax\else\textmd{ : \itshape #1}\fi}}
\newtcolorbox{methode}[1][]{boxbase,colback=methcol!7,colframe=methcol,sharp corners,
  title={\textsc{Méthode}\ifx\relax#1\relax\relax\else\textmd{ --- \itshape #1}\fi}}
\newtcolorbox{exemplebox}[1][]{boxbase,colback=excol!5,colframe=excol,
  title={\textsc{Exemple}\ifx\relax#1\relax\relax\else\textmd{ : \itshape #1}\fi}}
\newtcolorbox{piege}[1][]{boxbase,colback=attncol!6,colframe=attncol,
  title={\textsc{Piège}\ifx\relax#1\relax\relax\else\textmd{ : \itshape #1}\fi}}
\newtcolorbox{remarque}[1][]{boxbase,colback=black!4,colframe=black!45,
  title={\textsc{Remarque}\ifx\relax#1\relax\relax\else\textmd{ : \itshape #1}\fi}}
% --- boîtes EXERCICE / CORRIGÉ (noms EXACTS, auto-comptées) ---
\newtcolorbox[auto counter]{exo}[1][]{boxbase,colback=primary!4,colframe=primary,
  title={\textsc{Exercice}~\thetcbcounter\ifx\relax#1\relax\relax\else\textmd{ --- \itshape #1}\fi}}
\newtcolorbox[auto counter]{corrige}[1][]{boxbase,colback=excol!4,colframe=excol,
  title={\textsc{Corrigé}~\thetcbcounter\ifx\relax#1\relax\relax\else\textmd{ --- \itshape #1}\fi}}
\newcommand{\vv}[1]{\vec{#1}}\newcommand{\ds}{\displaystyle}
\newcommand{\R}{\mathbb{R}}\newcommand{\N}{\mathbb{N}}\newcommand{\Z}{\mathbb{Z}}
% (un \usepackage{fancyhdr}+en-tête est possible ; il est ignoré côté web)
% =========================================================================
\begin{document}

\begin{center}
{\large\bfseries\color{excol} Thème 1 — Corrigé Moyen}\\[-2pt]
{\color{excol}\rule{5cm}{1pt}}
\end{center}

\begin{corrige}[une sonde à ultrasons]
\begin{enumerate}[label=\alph*)]
  \item $f=\SI{2}{\kilo\hertz}=\SI{2000}{\hertz}$, donc
        $T=\dfrac{1}{f}=\dfrac{1}{2000}=\SI{5e-4}{\second}=\SI{0.5}{\milli\second}.$
  \item $\lambda=\dfrac{v}{f}=\dfrac{\SI{1500}{\metre\per\second}}{\SI{2000}{\hertz}}
        =\SI{0.75}{\metre}.$
\end{enumerate}
\end{corrige}

\begin{corrige}[du graphe spatial à la période]
Le graphe est en fonction de $x$ : on y lit une longueur d'onde. Ensuite $\lambda=vT$.
\begin{enumerate}[label=\alph*)]
  \item Le motif se répète tous les \SI{3}{\metre} : $\lambda=\SI{3}{\metre}.$
  \item $\lambda=vT\ \Rightarrow\ T=\dfrac{\lambda}{v}=\dfrac{\SI{3}{\metre}}{\SI{6}{\metre\per\second}}
        =\SI{0.5}{\second}.$
  \item $f=\dfrac{1}{T}=\dfrac{1}{0{,}5}=\SI{2}{\hertz}.$
\end{enumerate}
\end{corrige}

\begin{corrige}[le retard d'un point éloigné]
Le retard, c'est le temps que met la perturbation pour parcourir la distance jusqu'à P :
$\tau=x/v$.
\begin{enumerate}[label=\alph*)]
  \item $\tau=\dfrac{x}{v}=\dfrac{\SI{12}{\metre}}{\SI{3}{\metre\per\second}}=\SI{4}{\second}.$
  \item La source démarre à $t=0$ et P reproduit son mouvement avec le retard $\tau$ : P
        commence donc à bouger à $t=\tau=\SI{4}{\second}$.
\end{enumerate}
\end{corrige}

\begin{corrige}[le son change de milieu]
La fréquence est imposée par la source ; seules $v$ et $\lambda$ dépendent du milieu.
\begin{enumerate}[label=\alph*)]
  \item $\lambda_{\text{air}}=\dfrac{v_{\text{air}}}{f}
        =\dfrac{\SI{340}{\metre\per\second}}{\SI{680}{\hertz}}=\SI{0.5}{\metre}.$
  \item \textbf{Non}, $f$ ne change pas : elle est fixée par la source (le diapason). On garde
        $f=\SI{680}{\hertz}$.
  \item $\lambda_{\text{eau}}=\dfrac{v_{\text{eau}}}{f}
        =\dfrac{\SI{1360}{\metre\per\second}}{\SI{680}{\hertz}}=\SI{2}{\metre}$ :
        la longueur d'onde est 4 fois plus grande dans l'eau (car $v$ y est 4 fois plus grand).
\end{enumerate}
\end{corrige}

\begin{corrige}[la bouée qui oscille]
\begin{enumerate}[label=\alph*)]
  \item 20 oscillations en \SI{10}{\second} :
        $T=\dfrac{\SI{10}{\second}}{20}=\SI{0.5}{\second}$, d'où
        $f=\dfrac{1}{T}=\SI{2}{\hertz}$ (ou directement $f=20/10=\SI{2}{\hertz}$).
  \item La distance entre deux crêtes voisines \emph{est} la longueur d'onde :
        $\lambda=\SI{3}{\metre}.$
  \item $v=\lambda f=\SI{3}{\metre}\times\SI{2}{\hertz}=\SI{6}{\metre\per\second}.$
\end{enumerate}
\end{corrige}

\end{document}
